% !TeX root = ../papers/nc-p3.tex
\section{十八次の滑らかな二次核曲面の二重被覆の排除}
\label{sec:cofactor-eighteen-double-quadric}
\begin{theorem}\label{thm:cofactor-eighteen-no-smooth-quadric}
$d=18$ では滑らかな二次曲面は核像にならない。
\end{theorem}
\begin{proof}[二つの母線次数]
前定理によって二重被覆だけを扱う。$Q$ を滑らかな核二次曲面、$T$ を入力八次の関数体での
$Q$ の正規化とする。$T\to Q$ は有限平坦二重被覆である（正規曲面はCohen--Macaulayで底は滑らか）。
跡分解により被覆代数は $\mathcal O_Q\oplus\mathcal O_Q(-a,-b)$、
積は次数 $(2a,2b)$ の被約な分岐式による（平方因子があれば正規でない）。
核グラフを解消した $S$ 上で二母線の引戻しを $A,B$ とし、$L=A+B$。
$A^2=B^2=0$, $AB=2$, $H^2=8$, $HL=6$ である。$A,B$ の包への直交射影で
$(HA)(HB)\ge8$、順序を除き $(HA,HB)=(2,4),(3,3)$ である。

$(2,4)$ は等号なので $H$ は $A,B$ の正の線形結合に数値同値である。
従って入力正規化は $T$ に一致し、入力射 $f:T\to D$ は有限である（両収縮の共通解消で潰す曲線が一致）。
対称余核 $M$ と有限双有理な評価 $M\subset f_*L$ の跡双対から
$2L-K_T-2H$ の非零切断を得る。詳しくは定理\ref{thm:kernel-curve-scroll}と同じ
$f_*(\omega_T L^{-1})\subset M(-2H)\subset f_*L(-2H)$ の合成が
分数体の乗法であり $\mathcal O_T$ 線形である。入力次数四の母線逆像が一般に連結なら
滑らかなその曲線上でこの切断の次数は $4-(2g-2)-8<0$ となる。
従って一般逆像は非連結で、分岐式はその母線方向の次数零である。
二重被覆は $B_0\times\mathbb P^1\to\mathbb P^1\times\mathbb P^1$、
第一因子での二重被覆になる。もう一方の一般母線逆像 $B_0$ は入力次数二で、
入力写像は一般曲線上で双有理なので滑らかな有理曲線である。
よって $T=\mathbb P^1\times\mathbb P^1$, $H=(2,2)$, $L=(2,1)$。
上の切断の零因子は類 $(2,0)$ で空でない。その各成分は第一座標が固定した母線である。
その入力像は次数二以下なので整な十八次の階数二曲線の成分ではなく、一般点は階数三である。
一方、第二座標が一般なら $f$ は第一因子の曲線を滑らかな円錐曲線として埋める。
従って第二座標と $f$ の対はその開集合で閉埋込みである（各繊維の相対二次完全系）。
階数三の入力点の近くでは核座標は入力から正則に復元され、他の第二座標を持つ逆像もない。
従って $f$ は上記零因子の一般点の像で同型となり、そこで $M=f_*L$ とその双対も同型である。これは矛盾。
\end{proof}

\begin{proof}[三次の母線と底の種数]
$(3,3)$ では両方の一般母線逆像が連結である。非連結なら二成分のモノドロミーは推移的で、
全入力次数は偶数になる。一般逆像は滑らかで入力では三次曲線に双有理なので種数は零か一である。
一般入力切断の種数は三以下なので $K_S H\le-4$。
不正則数が零なら種数一の一般母線逆像 $F$ は存在しない。実際随伴完全列で
$H^0(K_S+F)\to H^0(K_F)$ は全射、従って $(K_S+F)H\ge0$ となり $FH=3$ に反する。

不正則数が正の場合も排除する。$K_SH<0$ と $H$ のnef性から小平次元は負、
従って曲面は線織である。両母線の一般逆像は線織面の底曲線を支配する
（支配しない動く曲線は有理母線で、この母線族自身ならもう一方の交差する族が底を支配し、
その有理母線族のパラメータ $\mathbb P^1$ が正種数の底を支配することになる）。
ゆえに底は種数一、両方とも種数一である。線織面の一般有理母線を $F_0$、$k=HF_0$ とする。
射影直線束モデルに入力系を移し固定部を除いたとき、基点重複度 $m_i$ は $k$ 以下である
（一般の系の一員と基点を通る母線を交わし、近接重複度は増えない）。
種数一の底上の直線束では、相対次数 $k$ の因子 $D_0$ は $kK D_0=-D_0^2$ を満たすので
\[
 K_S H=-\frac8k+\sum_i m_i(1-m_i/k)\ \ge -\frac8k .
\]
従って $k=1,2$。$k=1$ では全非零重複度は $k$、$k=2$ でも $K_SH\le-4$ から同様である。
初等変換により入力系に基点のない射影直線束モデルが得られる。
実際重複度 $k$ の点を吹けばそこの母線の固有変換と一般切断は交わらない。
これを潰せば新基点は生じず、元の一点分だけ基点が減るので反復できる。

このモデルで母線逆像の二鉛筆の像の一方を $U$ とし、$u=UF_0\ge1$。
$HU=3$, $H^2=8$、直線束の二次元数値格子で
$U^2=u(6k-8u)/k^2$ である。移動系ゆえ非負で、$k=1$ は不可能、$k=2$ なら $u=1$, $U^2=1$。
二鉛筆とも $HF_0$ と $HU$ によって同じ数値類を持ち、相互交点も一である。
元の二鉛筆の交点二は共通解消で基点を除くとこれ以下になるはずなので矛盾。
\end{proof}

\begin{proof}[次数四のdel Pezzo二重被覆]
従って両方の一般母線逆像は有理曲線で、分岐次数は $(2,2)$ である。
$T$ は有理曲面である（$\mathbb P^1$ 上の有理曲線束に対するTsenの定理）。
被覆代数から $\chi(\mathcal O_T)=1$、有理な解消との比較で全て有理特異点であり、
Gorensteinなので有理二重点である。最小解消 $S_0$ では $K_{S_0}=-L$。
逆入力系の固定部を除いた類 $H_0$ は $H_0 L=6$ で $8\le H_0^2\le9$。
随伴の偶奇から $H_0^2=8$ で基点はなく、以下 $S=S_0$, $H=H_0$ とする。一般入力切断の種数は二である。

$P=H-L=H+K_S$ は $P^2=0$, $HP=2$, $LP=2$ を満たす。
これは基点なしの有理曲線鉛筆である。実際有理性と随伴で $|P|$ の切断空間は
種数二の一般入力切断の標準切断空間に同型である。従って固定部 $Z$ は $HZ=0$。
移動部 $U$ は $HU=2$ とHodge指数により $0\le U^2\le1/2$、よって自由な鉛筆である。
一般繊維は連結である（Stein分解の底は有理で、完全な鉛筆の次数は一）。
$LU>0$ なので随伴から $LU=2$。従って $LZ=0$, $PZ=0$ となり $U^2=P^2$ から $Z^2=0$。
$H$ に直交する有効な $Z$ は負定値性で零となる。特に $H=L+P$ より入力が潰す曲線は全て $L$ も潰す。

入力射が有限な開集合では最初の証明の跡双対の合成から $2L-K_S-2H=3L-2H$ の有理切断が得られ、
そこで零因子は有効である。取り除いた入力点上の曲線は $L$ 次数零であり、極があってもその上に限られる。
$(3L-2H)L=0$ なので零因子の全成分も $L$ 次数零である。
従ってこれらの曲線の入力像以外では余次元一で $M=f_*L$ である
（跡の合成が各正規化枝で単元ならその中間層も押出し全体になる）。

各特異核被覆点 $z\in T$ の核像 $y\in Q$ で、$y$ を通る一般超平面の全引戻しの固定曲線部分を $Z_y$ とする。
これは $z$ 上の全例外曲線を含み $LZ_y=0$ である。$y$ を通る一般二平面は $Q$ 上の別の一点でも交わり、
その逆像は二点なので $(L-Z_y)^2\ge2$。ゆえに $-2\le Z_y^2<0$。
$H-3L/2$ の自乗は $-1$ なのでHodge指数から $0\le HZ_y\le\sqrt2$。
一般入力平面の滑らかな切断 $B_H$ がこの例外集合に会うなら、一般の一成分に横断的にただ一点で会い、
その固定重複度は一である。従って $B_H\to Q$ はその点で微分非零、かつ $y$ の逆像はそれ一つだけである
（被覆特異点は分岐上で、二重被覆のそこでの台は一点）。各該当成分の $H$ 次数は一だから入力像は直線である。

一般入力平面で評価 $M_H\subset f_*\mathcal O_{B_H}(L)$ を比べる。
上記の直線以外で余次元一の評価は既に同型で、一般平面は孤立例外を避ける。
各直線の一般点は整な十八次の階数二曲線の外で階数三である。核座標はそこで入力から正則に復元できる。
$B_H\to Q$ が上記の一点で非分岐かつこの核座標を持つ他の逆像がないので、
その入力点での正規化繊維は一点で長さ一である。ここでも評価は全射となる。
従って一般平面上で二層は等しいが、右辺のEuler数は次数六、種数二から五であり、左辺は四である。
この矛盾で二重被覆も排除できた。
\end{proof}
\path{computations/verify_cofactor_eighteen_double_quadric.py} は母線次数、種数一線織モデルの交点格子、
および最小解消上の直交射影を有理数で検算する。
